\begin{figure}[htbp]
\centering
\begin{tikzpicture}[
 box/.style={draw=black!70,rounded corners=1pt,align=center,
             font=\small,inner sep=5pt,text width=40mm},
 arrow/.style={-{Stealth[length=2mm]},thick}]
\node[box] (memory) at (0,0) {Carga y memorias\\\((q,m_0,m_1)\)};
\node[box] (register) at (0,-1.65) {Registro reconstruido\\\(K=(q/12)\mathbf1+Lm\)};
\node[box] (weights) at (-3.8,-3.55) {Partición ordenada\\\(d_j=K_j/q\)};
\node[box] (direction) at (3.8,-3.55) {Dirección de suma nula\\\(u=P_3P_{11}K\)};
\node[box] (geometry) at (-4.6,-6) {Medición de frontera\\y forma canónica\\\(C(d),\ \Omega_{P_K}\)};
\node[box] (energy) at (0,-6) {Energía nodal\\y detalles\\\(L_d,\ \mathcal E(z)\)};
\node[box] (lattice) at (4.6,-6) {Deformación reticular\\y red dual\\\(\Lambda_s^*=\Lambda_{-s}\)};
\draw[arrow] (memory)--(register);
\draw[arrow] (register)--(weights);
\draw[arrow] (register)--(direction);
\draw[arrow] (weights)--(geometry);
\draw[arrow] (weights)--(energy);
\draw[arrow] (direction)--(lattice);
\draw[arrow] (direction) to[bend left=12] node[right,font=\scriptsize,align=left,fill=white,inner sep=1pt]{flujo\\heredado} (energy);
\end{tikzpicture}
\caption{Realizaciones posteriores del registro recuperable. La rama
intervalar determina los menores de caminos, la forma canónica del
polígono y el laplaciano nodal. La dirección algebraica determina el
flujo de intervalos y, junto con la realización reticular marcada, la
colección dual de veinticuatro dimensiones. Cada flecha conserva el dominio
y las hipótesis de su construcción; los codominios inferiores son
distintos.}
\label{fig:multimorphology}
\end{figure}
